\section{Top orders divisible by four}\label{sec:four}

We keep the setting of \Cref{sec:even}: $M$ is even and prime to $3$, $A$ is
a sextuple at level $M$ that contains no pair, is not split and generates
$\ZZ/M$, $e^{*}$ is a maximal odd part of $A$ and $2^{J}e^{*}$ is its largest
order. Here $J\ge2$, so $4\mid M$. Replacing $A$ by $-A$ does not change the
problem (\Cref{lem:algtools}(c) with $t=-1$).

\subsection{The case \texorpdfstring{$J=2$}{J=2}: two relations}
For $v\in U^{*}=(\ZZ/e^{*})^{\times}$ let
\begin{gather*}
  x^{+}_{v},\ x^{-}_{v}=\tfrac{M}{4e^{*}}\,u\quad(u\equiv1,\ \text{resp.}\
  u\equiv3\bmod4,\ u\equiv v\bmod e^{*}),\\
  y_{v}=\tfrac{M}{2e^{*}}\,u\quad(u\ \text{odd},\ u\equiv v\bmod e^{*}),
\end{gather*}
and $z_{v}=(M/e^{*})v$, of orders $4e^{*}$, $2e^{*}$ and $e^{*}$. Then
\[
  -x^{+}_{v}=x^{-}_{-v},\quad x^{-}_{v}=x^{+}_{v}+\tfrac M2,\quad
  \pm2x^{\pm}_{v}=y_{\pm v},\quad y_{v}+\tfrac M2=z_{v/2},\quad
  z_{v}+\tfrac M2=y_{2v},\quad -2y_{v}=z_{-v},
\]
where $\pm2x^{\pm}_{v}=y_{\pm v}$ means $2x^{+}_{v}=2x^{-}_{v}=y_{v}$ and
$-2x^{\pm}_{v}=y_{-v}$. By the maximality of $e^{*}$ and of $J$, these are
the elements of odd part $e^{*}$; the other elements are \emph{low}. Let
$N_{x}$, $N_{y}$, $N_{z}$ be the numbers of elements of $A$ of orders $4e^{*}$,
$2e^{*}$, $e^{*}$, and put
\begin{gather*}
  P(v)=\mult_{A}(x^{+}_{v})-\mult_{A}(x^{-}_{-v}),\\
  Y(v)=\mult_{A}(y_{v})-\mult_{A}(y_{-v}),\qquad
  Z(v)=\mult_{A}(z_{v})-\mult_{A}(z_{-v}).
\end{gather*}
$Y$ and $Z$ are odd, and as $A$ contains no pair, $\sum_{v}|P(v)|=N_{x}$,
$\|Y\|=2N_{y}$ and $\|Z\|=2N_{z}$. With $\hat f(\chi)=\sum_{v}f(v)\bar
\chi(v)$ we have $\widehat{f(2\,\cdot)}=\chi(2)\hat f$.

\begin{lemma}\label{lem:R12}
\begin{enumerate}[label=\textup{(R\arabic*)}]
\item $\hat Z(\chi)+(1-\chi(2))\bigl(\hat Y(\chi)+\hat P(\chi)\bigr)=0$ for
every odd primitive character $\chi$ modulo $e^{*}$;
\item $\hat P(\psi)=0$ for every even primitive character $\psi$ modulo
$e^{*}$.
\end{enumerate}
Consequently $D(v)=P(v)+P(-v)$ and $G=Z+Y-Y(2\,\cdot)+P-P(2\,\cdot)$ are split
along $(K_{p})_{p\mid e^{*}}$ on every coset of $H^{*}$, and $e^{*}>1$.
\end{lemma}

\begin{proof}
(R1) In \eqref{eq:charrel} with $c=e^{*}$, the elements whose order is
divisible by $e^{*}$ are the $x^{\pm}_{v}$, $y_{v}$, $z_{v}$. The element
$z_{v}$ contributes $\bar\chi(v)/\varphi(e^{*})$, the element $y_{v}$
contributes $\bar\chi(v)(1-\chi(2))/\varphi(e^{*})$, and $x^{\pm}_{v}$
contributes $\bar\chi(v)(1-\chi(2))/(2\varphi(e^{*}))$. Since $\chi$ is odd,
$\sum_{v}\mult_{A}(z_{v})\bar\chi(v)=\hat Z(\chi)/2$, likewise for $Y$, and
$\sum_{v}(\mult_{A}(x^{+}_{v})+\mult_{A}(x^{-}_{v}))\bar\chi(v)=\hat
P(\chi)$.

(R2) A primitive character of conductor $4e^{*}$ is $\chi_{4}\psi$ with
$\psi$ primitive modulo $e^{*}$, and it is odd exactly when $\psi$ is even.
Only the elements of order $4e^{*}$ enter \eqref{eq:charrel}, and
$\overline{\chi_{4}\psi}(u)=\pm\bar\psi(v)$ for $x^{\pm}_{v}$; so
$\sum_{v}(\mult_{A}(x^{+}_{v})-\mult_{A}(x^{-}_{v}))\bar\psi(v)=0$, which is
$\hat P(\psi)=0$ as $\psi$ is even.

For a primitive $\chi$: if $\chi$ is odd, $\hat D(\chi)=0$ as $D$ is even, and
$\hat G(\chi)=\hat Z+(1-\chi(2))(\hat Y+\hat P)=0$ by (R1); if $\chi$ is even,
$\hat D(\chi)=2\hat P(\chi)=0$, and $\hat Z=\hat Y=0$, so $\hat
G(\chi)=(1-\chi(2))\hat P(\chi)=0$ by (R2). The proof of \Cref{lem:splitF}
gives the splitting. If $e^{*}=1$, all odd parts are $1$, $M=4$, and (R2) at
the conductor $4$ reads $\mult_{A}(M/4)=\mult_{A}(3M/4)$, so $N_{x}=0$ as $A$
contains no pair; but $J=2$.
\end{proof}

So $e^{*}\ge5$ is prime to $6$, and the subgroups $K_{p}$ satisfy the
hypotheses of \Cref{lem:S} and \Cref{lem:grid}.

\begin{lemma}\label{lem:fourmoves}
Each of the following is a move that leaves the elements of order larger
than $4e^{*}$ in place and either lowers the number of elements of order
$4e^{*}$, or keeps it and lowers the number of elements of order $2e^{*}$;
so it lowers the profile, unless at most four elements remain.
\begin{enumerate}[label=\textup{(\alph*)}]
\item $x^{+}_{v},x^{-}_{v}\in A$: with $x=x^{+}_{v}$ and
$S=\{x,x+M/2\}\subseteq T(x)$, $A''=A-x^{+}_{v}-x^{-}_{v}+\{y_{v},M/2\}$.
And $y_{2v},z_{v}\in A$: with $S=\{z_{v},z_{v}+M/2\}\subseteq T(z_{v})$,
$A''=A-y_{2v}-z_{v}+\{z_{2v},M/2\}$.
\item $x=x^{\pm}_{v}$ and $y_{-v}=-2x$ in $A$: with $S=\{x,-2x\}$,
$A''=A-x-y_{-v}+\{-x+M/2,\,M/2\}$. And $y_{v},z_{-v}\in A$: with
$S=\{y_{v},-2y_{v}\}\subseteq T(y_{v})$, $A''=A-y_{v}-z_{-v}+\{z_{-v/2},M/2\}$.
\item $y_{v},y_{2v}\in A$: with $\sigma=T(y_{2v})+T(y_{v})$ and
$S=\{y_{2v},y_{v}\}$, $A''=A-y_{v}-y_{2v}+\{z_{2v},z_{-v/2}\}$.
\item For a prime $q\mid e^{*}$ and a line $vK_{q}$: if $A$ contains
$(q+1)/2$ distinct elements of one of the four kinds $x^{+}_{w}$, $x^{-}_{w}$,
$y_{w}$, or $z_{w}$ with $e^{*}\ne q$, for $w\in vK_{q}$, then $A$ has a move.
\end{enumerate}
\end{lemma}

\begin{proof}
(a)--(c) are instances of \Cref{def:move}; in (c),
$-(\sigma-S)=\{z_{-v},z_{2v},M/2,z_{-v/2},z_{v},M/2\}$, from which the pairs
$\{z_{v},z_{-v}\}$ and $\{M/2,M/2\}$ are removed. In (b) the added element
$-x+M/2$ has order $4e^{*}$, and the element $y_{-v}$ of order $2e^{*}$ is
removed. For (d), as in \Cref{lem:topfibres}, the elements of one kind with
$w\in vK_{q}$ differ by elements of $(M/q)\ZZ$, so they lie in one fibre of
$\Phi_{q}$, and \Cref{lem:movecrit} applies, since $q$ differs from their
order.
\end{proof}

From now on $A$ has none of these moves.

\begin{lemma}\label{lem:foursigns}
\begin{enumerate}[label=\textup{(\roman*)}]
\item $P(v)$ and $P(-v)$ never have opposite signs. So
$|D(v)|=|P(v)|+|P(-v)|$, $\|D\|=2N_{x}\le12$, $D\ne0$, and where $P(v)\ne0$
it has the sign of $D(v)$.
\item $Z(w)$, $Y(w)$ and $-Y(2w)$ never have opposite signs. So
$|Z+Y-Y(2\,\cdot)|=|Z|+|Y|+|Y(2\,\cdot)|$.
\end{enumerate}
\end{lemma}

\begin{proof}
(i) $P(v)\ge1$ and $P(-v)\le-1$ mean $x^{+}_{v},x^{-}_{v}\in A$, and
$P(v)\le-1$ and $P(-v)\ge1$ mean $x^{-}_{-v},x^{+}_{-v}\in A$; both are
\Cref{lem:fourmoves}(a). And $N_{x}\ge1$ as $J=2$. (ii) Opposite signs of
$Z(w)$ and $Y(w)$ give $z_{\mp w}$ and $y_{\pm w}$ in $A$, which is (b); of
$Z(w)$ and $-Y(2w)$ give $z_{\pm w}$ and $y_{\pm2w}$, which is (a); of $Y(w)$
and $-Y(2w)$ give $y_{\pm w}$ and $y_{\pm2w}$, which is (c).
\end{proof}

\begin{lemma}\label{lem:evencore}
Up to replacing $A$ by $-A$, $D$ has one of three shapes:
\begin{enumerate}[label=\textup{(\Roman*)}]
\item $D=\mathbf 1_{L}+\mathbf 1_{-L}$ with $L\ne-L$ a line along $K_{q}$,
$|K_{q}|\le6$;
\item $e^{*}=p$ is prime and $D$ is a positive constant;
\item $e^{*}=35$ and $N_{x}=6$.
\end{enumerate}
\end{lemma}

\begin{proof}
Nonzero split functions of mass at most $8$ have mass at least $4$
(\Cref{lem:S}). Let $C$ be a coset of $H^{*}$ on which $D\ne0$.

\emph{Case (a): $-1\notin H^{*}$.} The cosets $C$ and $-C$ are distinct with
the same mass, so $D$ has mass at most $6$ on $C$, and by \Cref{lem:grid}
$D=\sigma\mathbf 1_{L}$ on $C$, with $L$ a line along some $K_{q}$ with
$|K_{q}|\le6$, and $D=\sigma\mathbf 1_{-L}$ on $-C$. These use mass
$2|L|\ge8$, and the other cosets, which come in pairs, carry nothing: shape I.

\emph{Case (b1): $-1\in H^{*}$ and $e^{*}=p$ is prime.} Then
$H^{*}=K_{p}=U^{*}$, and $D$ is constant by (S4); its value is nonzero, and
we make it positive: shape II. (If $e^{*}=p^{k}$ with $k\ge2$, then
$H^{*}=K_{p}$ consists of units congruent to $1$ modulo $p^{k-1}$, and
$-1\notin H^{*}$.)

\emph{Case (b2): $-1\in H^{*}$ and $e^{*}$ has at least two prime factors.}
Let $K=K_{p}$ have the largest order, $H'$ the product of the other $K_{q}$,
and write $-1=\varepsilon\varepsilon'$ with $\varepsilon\in K$,
$\varepsilon'\in H'$; by \Cref{lem:Kp}(b), $\varepsilon\ne1\ne\varepsilon'$.
For the slices $S_{c}$ of $C$, $c\in K$, multiplication by
$\varepsilon\varepsilon'=-1$ maps $S_{c}$ onto $S_{c\varepsilon}$, and as
$D$ is even, $\mu(c\varepsilon)=\mu(c)$. A coset $C'\ne C$ carrying a nonzero
part of $D$ satisfies $C'=-C'$, and an even nonzero split function of mass $4$
on it would be $\pm\mathbf 1_{L''}$ with $L''=-L''$, that is, $-1\in K_{q}$,
which \Cref{lem:Kp}(b) excludes; so other cosets carry mass at least $6$, or
nothing.

(i) If some slice vanishes, $D$ is split along $(K_{q})_{q\ne p}$ on every
slice by (S2). A nonzero slice has $\mu(c)=\mu(c\varepsilon)\le6$, and by
\Cref{lem:grid} $D=\sigma$ on a line $L$ along some $K_{q}$, $q\ne p$, with
$|K_{q}|\le6$, inside $S_{c}$ and $0$ elsewhere on $S_{c}$; on
$S_{c\varepsilon}$, $D=\sigma$ on $-L$. This uses mass $2|L|\ge8$, so nothing
else carries mass: shape I.

(ii) If all slices are nonzero and $\mu(c_{s})=1$, then $\mu(c_{s}
\varepsilon)=1$, and $z\mapsto D(z)-D(z\varepsilon)$ on $S_{c_{s}}$ is split
along $(K_{q})_{q\ne p}$ by (S1), of mass at most $2$, hence zero. For the
support point $z_{0}$ of $S_{c_{s}}$, the points $z_{0}\varepsilon$ and
$z_{0}\varepsilon\varepsilon'=-z_{0}$ of $S_{c_{s}\varepsilon}$ both carry
$D(z_{0})\ne0$; so they coincide and $\varepsilon'=1$, a contradiction.

(iii) If all $\mu(c)\ge2$, then $2|K|\le12$, so every $|K_{q}|\le6$ and the
primes of $e^{*}$ are $5$ and $7$, with $7\parallel e^{*}$. If $25\mid e^{*}$,
then $|K_{5}|=5$, $K=K_{7}$, and $\varepsilon'\in K_{5}$ satisfies
$\varepsilon'^{2}=1$ (the product being direct), so $\varepsilon'=1$, a
contradiction. So $e^{*}=35$, $|K|=6$, every $\mu(c)=2$, and the mass is
$12=2N_{x}$: shape III.
\end{proof}

\subsection{Shape I}

\begin{lemma}\label{lem:shapeI}
In shape \textup{I}, $A$ has a chain of moves that is direct or lowering.
\end{lemma}

\begin{proof}
By \Cref{lem:foursigns}(i), $P=\mathbf 1_{S}$ with $S\subseteq L\cup(-L)$
containing exactly one of $v,-v$ for each $v\in L$; so the elements of order
$4e^{*}$ are the $x^{+}_{w}$, $w\in S$, each once, and $N_{x}=|L|$. By
\Cref{lem:fourmoves}(d), $S\cap L$ and $S\cap(-L)$ have at most $(q-1)/2$
points each. If $25\mid e^{*}$ and $q=5$, then $|L|=5>4$; so $q\parallel e^{*}$,
$|L|=q-1$, and both sets have exactly $(q-1)/2$ points.

Let $x=x^{+}_{v_{0}}$ with $v_{0}\in S\cap L$. The fibre $\Phi_{q}(x)$
consists of the $x^{+}_{w}$, $w\in L$, and of one element $x_{0}$ of order
$4e^{*}/q$. The move $(\sigma_{q,x},S')$ with $S'=\{x^{+}_{w}:w\in S\cap L\}$
gives
\[
  A''=A-S'+\{-x^{+}_{w}:\ w\in L\setminus S\}+\{-x_{0},\,qx\}.
\]
For $w\in L\setminus S$ we have $-w\in S$ and $-x^{+}_{w}=x^{-}_{-w}=
x^{+}_{-w}+M/2$. So $A''$ contains the $(q-1)/2$ pairs $\{x^{+}_{u},x^{-}_{u}\}$,
$u\in S\cap(-L)$, and the moves of \Cref{lem:fourmoves}(a) replace each of them
by $\{y_{u},M/2\}$. The result,
\[
  A-\{x^{+}_{w}:\ w\in S\}+\{y_{u}:\ u\in S\cap(-L)\}+\tfrac{q-1}{2}\{M/2\}
  +\{-x_{0},\,qx\},
\]
has eight elements, among them at least two copies of $M/2$; after removing a
pair at most six remain, with no element of order $4e^{*}$ and none of larger
order that $A$ did not have.
\end{proof}

\subsection{Shape II}

\begin{lemma}\label{lem:shapeII}
In shape \textup{II}, either $A$ has a chain of moves that is direct or
lowering, or, up to multiplication by a unit congruent to $1$ modulo $4$, the
elements of $A$ of odd part $p$ are one of
\begin{enumerate}[label=\textup{(\alph*)}]
\item $\{x^{+}_{1},x^{+}_{2},z_{3}\}$, $p=5$;
\item $\{x^{+}_{1},x^{+}_{2},x^{+}_{4}\}$, $p=7$;
\item $\{x^{+}_{1},x^{+}_{2},x^{+}_{3},z_{4},z_{5}\}$, $p=7$;
\item $\{x^{+}_{1},x^{+}_{2},x^{+}_{4},x^{+}_{8},x^{+}_{5},z_{6}\}$, $p=11$.
\end{enumerate}
\end{lemma}

\begin{proof}
Let $D=c>0$. By \Cref{lem:foursigns}(i), $0\le P(v)\le c$ and
$P(v)+P(-v)=c$. The $w$ with $P(w)\ge1$ index distinct elements $x^{+}_{w}$ of
the line $U^{*}=K_{p}$, so there are at most $(p-1)/2$ of them by
\Cref{lem:fourmoves}(d), and at least one in each pair $\pm v$. Hence
$P=c\,\mathbf 1_{S}$ for a set $S$ of representatives of $U^{*}/\{\pm1\}$, and
$N_{x}=c(p-1)/2\le6$.

\emph{$c\ge2$.} Let $x=x^{+}_{v_{0}}$, $v_{0}\in S$. Now $\Phi_{p}(x)$ consists of
all the $x^{+}_{w}$ and one element $x_{0}$ of order $4$, and as in
\Cref{lem:shapeI} the move $\sigma_{p,x}$ followed by the moves
\Cref{lem:fourmoves}(a) on the pairs $\{x^{+}_{u},x^{-}_{u}\}$, $u\in S$, gives
\[
  A-2\{x^{+}_{u}:u\in S\}+\{y_{u}:u\in S\}+\tfrac{p-1}{2}\{M/2\}
  +\{-x_{0},\,px\},
\]
with eight elements, at least two copies of $M/2$ and fewer elements of order
$4p$ than $A$.

\emph{$c=1$.} $G$ is constant on $U^{*}$ by (S4), and $\sum_{v}G(v)=0$, as
$Z$, $Y$ are odd and $\sum P=\sum P(2\,\cdot)$; so $G=0$. With
$\psi=\mathbf 1_{S}-\mathbf 1_{-S}$ we have $P=(1+\psi)/2$, so
\[
  Z+Y-Y(2\,\cdot)=P(2\,\cdot)-P=\tfrac12\bigl(\psi(2\,\cdot)-\psi\bigr),
\]
which is $1$ on $\Phi=\{w\notin S:\ 2w\in S\}$, $-1$ on $-\Phi$ and $0$
elsewhere. If $y_{v}\in A$, then by \Cref{lem:foursigns}(ii) the left side is
at least $1$ at $v$, so $v\notin S$, $-v\in S$, and $y_{v}=-2x^{+}_{-v}$ with
$x^{+}_{-v}\in A$, which is \Cref{lem:fourmoves}(b). So $N_{y}=0$,
$Z=\mathbf 1_{\Phi}-\mathbf 1_{-\Phi}$, and the elements of order $p$ are the
$z_{w}$, $w\in\Phi$, each once. Along the cycles of $w\mapsto2w$ on
$(\ZZ/p)^{\times}$, $\psi$ changes sign $2|\Phi|$ times, and $(p-1)/2+|\Phi|
\le6$, so $p\le13$. The units of $\ZZ/M$ congruent to $1$ modulo $4$ act on
$U^{*}$ through all of $(\ZZ/p)^{\times}$ and preserve the kinds $x^{+}$,
$z$.
\begin{itemize}
\item $p=5$: one cycle of length $4$ with $-1=2^{2}$, so $\psi$ takes the
values $a,b,-a,-b$ along it and changes sign twice; $|\Phi|=1$, and up to
units $S=\{1,2\}$, $\Phi=\{3\}$. This is (a).
\item $p=7$: two cycles $\{1,2,4\}$ and $-\{1,2,4\}$, on which $\psi$ takes
opposite values; $|\Phi|$ is the number of sign changes on one cycle, $0$ or
$2$. Up to units, $S=\{1,2,4\}$, $\Phi=\emptyset$, which is (b), or
$S=\{1,2,3\}$, $\Phi=\{4,5\}$, which is (c).
\item $p=11$ and $p=13$: one cycle of length $p-1$ with $-1=2^{(p-1)/2}$, so
$\psi(2^{(p-1)/2}w)=-\psi(w)$ and the number of sign changes along a half of
the cycle is odd; $|\Phi|$ is odd. For $p=13$, $6+|\Phi|>6$. For $p=11$,
$|\Phi|=1$, $\psi=1$ on five consecutive powers of $2$, and up to units
$S=\{1,2,4,8,5\}$, $\Phi=\{6\}$. This is (d).\qedhere
\end{itemize}
\end{proof}

\subsection{Shape III}

\begin{lemma}\label{lem:shapeIII}
Shape \textup{III} does not occur.
\end{lemma}

\begin{proof}
Here $N_{x}=6$, so $Y=Z=0$, $G=P-P(2\,\cdot)$, and $\hat G(\chi)=(1-\chi(2))
\hat P(\chi)$ vanishes at every primitive character $\chi$ modulo $35$. Such a
character is $\chi_{5}\chi_{7}$ with both factors nontrivial; $\chi_{5}(2)$
has order $2$ or $4$, as $2$ generates $(\ZZ/5)^{\times}$, and $\chi_{7}(2)$
has order $1$ or $3$, as $2$ has order $3$ modulo $7$; so $\chi(2)\ne1$. Hence
$\hat P$ vanishes at every primitive character, $P$ is split on
$U^{*}=H^{*}$, nonzero, of mass $6$, and by \Cref{lem:grid} $P=\sigma\mathbf
1_{L}$ for a line $L$ along $K_{7}$. Then $A$ contains six distinct elements
$x^{+}_{w}$, $w\in L$, or six distinct elements $x^{-}_{-w}$, $w\in L$, on one
line along $K_{7}$, which \Cref{lem:fourmoves}(d) excludes.
\end{proof}

\subsection{The case \texorpdfstring{$J\ge3$}{J>=3}}

\begin{lemma}\label{lem:Jthree}
If $J\ge3$, then $A$ has a chain of moves that is direct or lowering.
\end{lemma}

\begin{proof}
The elements of order $e=2^{J}e^{*}$ are
$x_{t,v}=(M/e)u$ with $u\equiv t$ modulo $2^{J}$ and $u\equiv v$ modulo
$e^{*}$, for $(t,v)\in(\ZZ/2^{J})^{\times}\times U^{*}$. Put
$\tilde F(t,v)=\mult_{A}(x_{t,v})-\mult_{A}(x_{-t,-v})$. Since $e$ is a
maximal order, \eqref{eq:charrel} with $c=e$ involves only these elements, as
in \Cref{sec:odd}, so $\hat{\tilde F}$ vanishes at every primitive character
modulo $e$ (at the even ones because $\tilde F$ is odd). A character modulo
$e$ that is not primitive is trivial on the units congruent to $1$ modulo
$e/l$ for some prime $l\mid e$; for $l=2$ these form $\tilde K_{2}=\{1,\eta\}
\times\{1\}$ with $\eta=1+2^{J-1}$, and for $l=p\mid e^{*}$ they form
$\{1\}\times K_{p}$. So $\tilde F$ is split along $\tilde K_{2}$ and the
$K_{p}$. As $x_{t,v}+M/2=x_{t\eta,v}$, the pairs $\{x_{t,v},x_{t\eta,v}\}$ are
the pairs $\{x,x+M/2\}$.

Suppose $A$ contains no such pair, as otherwise the move $(T(x),\{x,x+M/2\})$
lowers the profile. By (S1), $g_{t}(v)=\tilde F(t,v)-\tilde F(t\eta,v)$ is
split along $(K_{p})_{p\mid e^{*}}$ on every coset of $H^{*}$. As $J\ge3$, the
orbit $O=\{t,t\eta,-t,-t\eta\}$ has four elements. Of the four elements
$x_{t,v}$, $x_{t\eta,v}$, $x_{-t,-v}$, $x_{-t\eta,-v}$, two opposite ones
form a pair and two adjacent ones a pair $\{x,x+M/2\}$, so $A$ contains copies
of at most one of them or of the diagonal $\{x_{t,v},x_{-t\eta,-v}\}$ or
$\{x_{t\eta,v},x_{-t,-v}\}$, and $|g_{t}(v)|$ is the number of elements of
$A$ among the four. Summing over $v$, the mass of $g_{t}$ is the number of
elements of $A$ of order $e$ whose first coordinate lies in $O$. A nonzero
split function has mass at least $4$, and the number of elements of order $e$
is at most $6$; so exactly one orbit $O$ is occupied, and for $t_{0}\in O$,
$g_{t_{0}}=\sigma\mathbf 1_{L}$ on one coset with $L$ a line along some
$K_{q}$, $|L|\le6$ (\Cref{lem:grid}). If $e^{*}=1$, the family $(K_{p})$ is
empty, $g_{t_{0}}=0$, and $A$ has no element of order $e$, a contradiction.

Replacing $t_{0}$ by $t_{0}\eta$ we may take $\sigma=1$. Then for $v\in L$
exactly one of $x_{t_{0},v}$ and $x_{-t_{0}\eta,-v}$ lies in $A$, once, and
$A$ has no other element of order $e$. The $x_{t_{0},v}$, $v\in L$, lie in one
fibre of $\Phi_{q}$, and so do the $x_{-t_{0}\eta,-v}$; so by
\Cref{lem:movecrit} each kind occurs at most $(q-1)/2$ times, $q\parallel
e^{*}$, $|L|=q-1$, and both occur $(q-1)/2$ times. The move $\sigma_{q,x}$ with
$x=x_{t_{0},v_{0}}\in A$ adds $-x_{t_{0},w}=x_{-t_{0},-w}$ for the $w\in L$
with $x_{t_{0},w}\notin A$, and these form pairs
$\{x_{-t_{0},-w},x_{-t_{0}\eta,-w}\}$ with elements of $A$. As in
\Cref{lem:shapeI}, the moves $T$ on these pairs leave at most six elements
after the removal of a pair of copies of $M/2$, none of order $e$, and no
element of larger order that $A$ did not have.
\end{proof}

\subsection{The global step}

\begin{proposition}\label{prop:globalfour}
Let $M$ be even and prime to $3$, with $M\notin\{70,110\}$. Every sextuple
$A$ at level $M$ that contains no pair, is not split and generates $\ZZ/M$ has
a chain of moves that succeeds.
\end{proposition}

\begin{proof}
Suppose not. By \Cref{lem:maxeven}, every maximal odd part $e^{*}$ of $A$ has
$J\ge1$; $J=1$ is excluded by \Cref{thm:twoparallel} and $J\ge3$ by
\Cref{lem:Jthree}. So every maximal odd part has $J=2$, and $4\mid M$. By
\Cref{lem:evencore,lem:shapeI,lem:shapeII,lem:shapeIII}, applied to $A$ or to
$-A$, the elements of odd part $e^{*}$ form, up to units and sign, one of the
configurations (a)--(d) of \Cref{lem:shapeII}, of sizes $3$, $3$, $5$, $6$.

If $A$ has two maximal odd parts, their configurations have three elements
each, with different primes: (a) for $5$ and (b) for $7$. Then $A$ lies in
the subgroup of order $140$, so $M=140$.

If $A$ has one maximal odd part $e^{*}=p\in\{5,7,11\}$, the low elements have
odd part $1$. Let $2^{K}$ be the largest order among them, and suppose
$K\ge3$. In \eqref{eq:charrel} at the conductor $2^{K}$ only the low elements
of order $2^{K}$ enter, as the elements of odd part $p$ have $2$-part at most
$4$; so the odd counting function of these elements on
$(\ZZ/2^{K})^{\times}$ has a vanishing transform at every primitive
character, and the other characters are trivial on $\{1,1+2^{K-1}\}$. Hence it
is invariant under multiplication by $1+2^{K-1}$, and an element $a\in A$ of
order $2^{K}$ comes with $(1+2^{K-1})a=a+M/2$; the move
$(T(a),\{a,a+M/2\})$ lowers the profile. So the low elements lie in
$\{M/4,M/2,3M/4\}$, $A$ lies in the subgroup of order $4p$, and
$M\in\{20,28,44\}$.

It remains to examine these levels. At $t=1$, balance says that the
representatives in $[1,M-1]$ of the elements of $A$ add up to $3M$, and we
may normalise $A$ by a unit and a sign.
\begin{itemize}
\item $M=20$, (a) $=\{1,17,12\}$: the three low elements lie in
$\{5,10,15\}$ and add up to $30$, so they are $\{5,10,15\}$ or
$\{10,10,10\}$, and $A$ contains a pair.
\item $M=28$, (b) $=\{1,9,25\}$: the low elements lie in $\{7,14,21\}$ and add
up to $49$, so they are $\{7,21,21\}$ or $\{14,14,21\}$, with a pair.
\item $M=28$, (c) $=\{1,9,17,16,20\}$: the low element is $21$. The moves
$(T(1),\{1\})$ and $(T(6),\{20,16,14\})$ turn
$A=\{1,9,16,17,20,21\}$ into $\{2,9,13,14,16,17,20,21\}$ and then into
$\{2,9,17\}+\{13,21,22\}$, a split sextuple; this chain succeeds.
\item $M=44$, (d) $=\{1,13,37,41,5,24\}$: the sum is $121\ne132$.
\item $M=140$, (a) for $5$ and (b) for $7$: up to units and sign, $A$ is
$\{7,119,84\}$ together with $\{5,45,125\}$ or with $\{55,75,115\}$, of sums
$385$ and $455$, not $420$.
\end{itemize}
So in every case $A$ contains a pair, is not balanced, or has a chain that
succeeds.
\end{proof}

\begin{theorem}\label{thm:primetothree}
The Hodge conjecture holds for $X^{4}_{M}$ when $3\nmid M$ and $110\nmid M$.
\end{theorem}

\begin{proof}
By induction on $M$. For $M$ odd this is \Cref{cor:odd}, and for $M=70$ it is
\Cref{cor:degrees}. Let $M$ be even, prime to $3$, not divisible by $110$
and different from $70$. The proper divisors of $M$ satisfy the same
conditions, so the Hodge conjecture holds for them, and
\Cref{prop:globalfour} with \Cref{prop:induction} gives it for $M$.
\end{proof}
